\documentclass[11pt]{article}
\usepackage[margin=1in]{geometry}
\usepackage{xcolor}
\usepackage{tcolorbox}
\usepackage{hyperref}
\usepackage{enumitem}
\setlist{noitemsep,topsep=0pt}
\usepackage{booktabs}
\usepackage{listings}
\usepackage{amsmath}
\usepackage{amssymb}
\usepackage{longtable}

% Define OSU orange color
\definecolor{osaorange}{RGB}{220,115,38}

% Configure hyperref colors
\hypersetup{
    colorlinks=true,
    linkcolor=blue,
    urlcolor=blue,
    citecolor=blue
}

% Configure R code listings
\lstset{
    language=R,
    basicstyle=\ttfamily\small,
    breaklines=true,
    frame=single,
    backgroundcolor=\color{gray!10},
    numbers=left,
    numberstyle=\tiny,
    keywordstyle=\color{blue},
    commentstyle=\color{green!60!black},
    stringstyle=\color{red}
}

\title{}
\author{}
\date{}

\begin{document}

% Orange header box with black outline
\begin{tcolorbox}[colback=osaorange,colframe=black,boxrule=2pt,arc=0pt,left=6pt,right=6pt,top=10pt,bottom=10pt]
\centering
\textcolor{white}{\textbf{\Large In-Class Worksheet}}\\
\textcolor{white}{\textbf{\large Rates, Standardization, and Advanced R Programming}}
\end{tcolorbox}

\vspace{8pt}

\noindent\textbf{Name:} \rule{8cm}{0.4pt} \hfill \textbf{Date:} \rule{4cm}{0.4pt}\\[0.5ex]
\textbf{Course:} H524 Introduction to Biostatistics\\
\textbf{Topic:} Rates and Standardization, Advanced R Programming, AI Integration

\vspace{8pt}

\section{Part 1: Review and Warm-up}

\subsection{Exercise 1.1: Quick Visualization Review}
Using the guinea pig data, work with a partner to create one boxplot comparing all three groups:

\begin{lstlisting}
# Recreate the data
control <- c(312, 345, 287, 398, 421, 356, 389, 334, 367, 401)
low_dose <- c(289, 267, 301, 234, 198, 245, 278, 213, 256, 298)
high_dose <- c(156, 189, 134, 167, 123, 145, 178, 201, 143, 134)

# Your code here:

\end{lstlisting}

\textbf{Quick Questions:}
\begin{itemize}
\item Which group has the most consistent survival times?
\item How do the medians compare across groups?
\end{itemize}

\section{Part 2: Introduction to Rates}

\subsection{Exercise 2.1: Understanding Rates in Public Health}

\textbf{Scenario:} You're analyzing COVID-19 data from two counties:

\textbf{County A:} 150 cases out of 10,000 residents\\
\textbf{County B:} 300 cases out of 25,000 residents

Calculate and compare:

\begin{lstlisting}
# County A
cases_A <- 150
population_A <- 10000
rate_A <- (cases_A / population_A) * 1000  # Rate per 1,000

# County B
cases_B <- 300
population_B <- 25000
rate_B <- (cases_B / population_B) * 1000  # Rate per 1,000

print(paste("County A rate per 1,000:", rate_A))
print(paste("County B rate per 1,000:", rate_B))
\end{lstlisting}

\textbf{Fill in your results:}
\begin{itemize}
\item County A raw cases: \rule{2cm}{0.4pt} out of \rule{2cm}{0.4pt}
\item County B raw cases: \rule{2cm}{0.4pt} out of \rule{2cm}{0.4pt}
\item County A rate per 1,000: \rule{2cm}{0.4pt}
\item County B rate per 1,000: \rule{2cm}{0.4pt}
\item Which county has the higher infection rate? \rule{3cm}{0.4pt}
\end{itemize}

\subsection{Exercise 2.2: Age-Standardized Rates}

\textbf{Extended Scenario:} Now consider age breakdown:

\textbf{County A:}
\begin{itemize}
\item Ages 0-50: 80 cases out of 7,000 residents
\item Ages 50+: 70 cases out of 3,000 residents
\end{itemize}

\textbf{County B:}
\begin{itemize}
\item Ages 0-50: 120 cases out of 15,000 residents
\item Ages 50+: 180 cases out of 10,000 residents
\end{itemize}

\begin{lstlisting}
# County A age-specific rates
county_A_young <- (80/7000) * 1000
county_A_old <- (70/3000) * 1000

# County B age-specific rates
county_B_young <- (120/15000) * 1000
county_B_old <- (180/10000) * 1000

print("Age-specific rates per 1,000:")
print(paste("County A (0-50):", round(county_A_young, 2)))
print(paste("County A (50+):", round(county_A_old, 2)))
print(paste("County B (0-50):", round(county_B_young, 2)))
print(paste("County B (50+):", round(county_B_old, 2)))
\end{lstlisting}

\textbf{Analysis Questions:}
\begin{enumerate}
\item Which age group has higher infection rates in both counties?
\item How does age structure affect the overall county comparison?
\item Why might age-standardization be important for fair comparisons?
\end{enumerate}

\section{Part 3: Advanced R Programming}

\subsection{Exercise 3.1: Creating Functions}

Based on our rate calculations, create a reusable function:

\begin{lstlisting}
# Template function - complete the missing parts
calculate_rate <- function(cases, population, per = 1000) {
  # Your code here
  # 1. Calculate the rate
  # 2. Add error checking for negative values
  # 3. Return formatted result

}

# Test your function
calculate_rate(150, 10000, 1000)
calculate_rate(300, 25000, 1000)
\end{lstlisting}

\subsection{Exercise 3.2: Working with Real Data Structures}

\begin{lstlisting}
# Create a more complex dataset
health_data <- data.frame(
  county = c("A", "A", "B", "B"),
  age_group = c("0-50", "50+", "0-50", "50+"),
  cases = c(80, 70, 120, 180),
  population = c(7000, 3000, 15000, 10000)
)

# Calculate rates for each row
health_data$rate_per_1000 <- (health_data$cases / health_data$population) * 1000

# Group calculations
library(dplyr)  # Load dplyr for data manipulation

# Calculate summary statistics by county
county_summary <- health_data %>%
  group_by(county) %>%
  summarise(
    total_cases = sum(cases),
    total_population = sum(population),
    crude_rate = (total_cases / total_population) * 1000
  )

print(county_summary)
\end{lstlisting}

\textbf{Interpretation Questions:}
\begin{enumerate}
\item What is the crude rate for each county?
\item How do these compare to your earlier calculations?
\item What additional information would help interpret these rates?
\end{enumerate}

\section{Part 4: Statistical Thinking and Study Design}

\subsection{Exercise 4.1: Connecting Rates to Biostatistics}

\textbf{Case Study Discussion:} A pharmaceutical company reports their new drug reduces heart attack rates from 4 per 1,000 patients to 2 per 1,000 patients.

\textbf{Questions for Discussion:}
\begin{enumerate}
\item What does this represent in terms of:
\begin{itemize}
\item Absolute risk reduction?
\item Relative risk reduction?
\end{itemize}
\end{enumerate}

\begin{lstlisting}
# Calculate the reductions
baseline_rate <- 4/1000
treatment_rate <- 2/1000

absolute_reduction <- baseline_rate - treatment_rate
relative_reduction <- (baseline_rate - treatment_rate) / baseline_rate

print(paste("Absolute risk reduction:", absolute_reduction))
print(paste("Relative risk reduction:", relative_reduction))
\end{lstlisting}

\begin{enumerate}
\setcounter{enumi}{1}
\item If you had 10,000 patients, how many heart attacks would be prevented?
\item What additional information would you want before recommending this drug?
\end{enumerate}

\subsection{Exercise 4.2: Bias and Confounding Preview}

Looking back at our county COVID comparison:
\begin{enumerate}
\item What factors besides county might affect infection rates?
\item How might differences in testing rates affect our rate calculations?
\item What is the population we can generalize these results to?
\end{enumerate}

\section{Part 5: Preparation for Future Topics}

\subsection{Exercise 5.1: Probability Preview}

When we cover probability theory, Based on today's work:

\begin{enumerate}
\item If County A's infection rate is 15 per 1,000, what's the probability that a randomly selected resident is infected?

\item How does this connect to what we've learned about populations and samples?
\end{enumerate}

\textbf{Bridge to Probability:}
\begin{itemize}
\item Rates can be thought of as probabilities
\item Population parameters vs. sample statistics
\item Uncertainty in our estimates
\end{itemize}

\section{Wrap-up Discussion}

\textbf{Key Concepts from Today:}
\begin{itemize}
\item Rates allow fair comparison between different-sized populations
\item Age standardization helps control for demographic differences
\item Functions in R make calculations reusable and less error-prone
\item Careful verification is essential for accurate statistical analysis
\end{itemize}

\textbf{Next Class Preview:} We'll begin with probability concepts - the foundation for all statistical inference.

\end{document}